% =====================================================================
% Submission abstract (LaTeX) — How Aggregation Can Conceal Composition: Aggregate
% Biocapacity and the Identifiability of Modelled Ecological-Capital Drawdown (Revision 35,
% two-land). Mathematical notation is kept as LaTeX (as required for the submitted
% manuscript).
% Data: National Footprint and Biocapacity Accounts (world, 2025 edition).
% =====================================================================
\begin{abstract}
Humanity's headline index of annual bioproductive supply --- total biocapacity --- is a single number built from heterogeneous land accounts: cropland, grazing land, forest products, fishing grounds, and built-up land. This paper asks whether that single number can conceal the composition of what is growing. It can, and the concealment is structural. From 1961 to 2022, global biocapacity grew about 23\%, but measured land-type accounts show that this growth came almost entirely from cropland, whose biocapacity grew 2.85$\times$, while the non-cropland book remained essentially flat at the endpoint. Over the same period, per-capita biocapacity fell by about half as population grew. A deliberately simple two-land model makes the mechanism precise: conversion from ecological-capital land to provisioning land can raise the measured total while the model's capital book falls; the drawdown is demand-driven and approaches a floor rather than showing slow precursors; and recovery requires an explicitly gated restoration flow. The deeper point is measurement. Because aggregate biocapacity combines areas, yields, and equivalence factors, the aggregate does not identify its own composition; a rising total may reflect broad growth, fast-provisioning expansion, or offsetting changes among components, and these cannot be separated from the total alone. Identifying composition requires independent land-cover and yield data, and global sustainability indicators should therefore report components alongside totals. These results establish an identifiability property of aggregate biocapacity and a measured composition pattern for 1961--2022; they are not a forecast and do not show a measured decline in the model's ecological-capital book.
\end{abstract}

\noindent\textbf{Keywords: } carrying capacity; biocapacity; ecological footprint; land-use conversion; composition illusion; identifiability; critical slowing down
